\chapter{الدوال المرجعية}\label{ch:g10:reffunc}

تظهر حفنة من الدوال البسيطة --- التآلفية، والمربع، والمقلوب، والجذر
التربيعي، والتكعيب --- في كل مكان، وحدها أو مركّبة. ومعرفة منحنياتها
وتغيراتها عن ظهر قلب تحوّل مسائل كثيرة إلى رسم سريع. يدرس هذا الفصل
كلًّا منها على حدة ويبرهن على تغيراتها بتقنية المقارنة في \cref{ch:g10:functions}.

\section{الدوال التآلفية}

\begin{definition}[الدالة التآلفية]\label{def:g10:reffunc:affine}
\emph{الدالة التآلفية}\index{دالة تآلفية} هي \omterm{def:g10:functions:function}{دالة} على الصورة
\[
f(x) = mx + p,
\]
حيث $m$ و $p$ عددان حقيقيان ثابتان. ومنحنيها مستقيم:
العدد $m$ هو \emph{الميل}\index{ميل} والعدد $p$ هو
\emph{الترتيب عند المبدأ}\index{الترتيب عند المبدأ} (فالمستقيم يقطع
المحور الشاقولي في $(0, p)$). وعندما $p = 0$ تكون \omterm{def:g10:functions:function}{الدالة}
\emph{خطية}: فالكتابة $f(x) = mx$ تعبّر عن التناسب.
\end{definition}

\begin{proposition}[الميل والتغيرات]\label{prop:g10:reffunc:affine}
لتكن $f(x) = mx + p$.
\begin{enumerate}
  \item من أجل أي دخلين مختلفين $u \neq v$:
        $m = \dfrac{f(v) - f(u)}{v - u}$ (\omterm{def:g10:reffunc:affine}{فالميل} هو تغير
        الخرج لكل وحدة تغير في الدخل).
  \item إذا كان $m > 0$، فإن $f$ \omterm{def:g10:functions:variations}{متزايدة} تمامًا على $\R$؛ وإذا كان $m < 0$،
        فهي \omterm{def:g10:functions:variations}{متناقصة} تمامًا؛ وإذا كان $m = 0$، فهي ثابتة.
\end{enumerate}
\end{proposition}

\begin{proof}
1. احسب $f(v) - f(u) = (mv + p) - (mu + p) = m(v - u)$، ثم اقسم على
$v - u \neq 0$.

2. إذا كان $u < v$، فإن $f(v) - f(u) = m(v-u)$ له إشارة $m$، لأن
$v - u > 0$: فإذا كان $m$ موجبًا كان $f(u) < f(v)$ (\omterm{def:g10:functions:variations}{متزايدة})، وإذا كان $m$ سالبًا
كان $f(u) > f(v)$ (\omterm{def:g10:functions:variations}{متناقصة}).
\end{proof}

\begin{omfigure}
\begin{tikzpicture}
\begin{axis}[omaxis,
  width=7.8cm, height=5.6cm,
  xmin=-3.4, xmax=3.4, ymin=-2.8, ymax=4.4,
  xtick={-2,-1,1,2}, ytick={-1,1,2,3}]
  \addplot[omDef, thick, domain=-3:2.6] {0.5*x + 1};
  \addplot[omThm, thick, domain=-2.2:3] {-x + 1};
  \node[omDef, font=\small] at (axis cs:2.5,2.7) {$y = \tfrac12 x + 1$};
  \node[omThm, font=\small] at (axis cs:-2.0,2.6) {$y = -x + 1$};
  \draw[dashed] (axis cs:1,1.5) -- (axis cs:2,1.5) -- (axis cs:2,2);
  \node[font=\footnotesize, below] at (axis cs:1.5,1.5) {$1$};
  \node[font=\footnotesize, right] at (axis cs:2,1.75) {$m$};
\end{axis}
\end{tikzpicture}

{\small دالتان تآلفيتان: \omterm{def:g10:functions:variations}{متزايدة} من أجل $m = \frac12 > 0$، \omterm{def:g10:functions:variations}{ومتناقصة}
من أجل $m = -1 < 0$. والانتقال وحدة واحدة نحو اليمين يغيّر $y$ بمقدار $m$.}
\end{omfigure}

\begin{example}\label{ex:g10:reffunc:affine}
أوجد \omterm{def:g10:reffunc:affine}{الدالة التآلفية} التي يمر منحنيها بالنقطتين $A(1, 3)$ و
$B(4, 9)$. \omterm{def:g10:reffunc:affine}{الميل} هو
\[
m = \frac{9 - 3}{4 - 1} = \frac{6}{3} = 2 .
\]
عندئذٍ $f(x) = 2x + p$، ويعطي $f(1) = 3$ أن $2 + p = 3$، إذن $p = 1$:
أي $f(x) = 2x + 1$. تحقق بالنقطة $B$: $f(4) = 9$.
\end{example}

\section{دالة المربع}

\begin{proposition}[دالة المربع]\label{prop:g10:reffunc:square}
\omterm{def:g10:functions:function}{الدالة} $f(x) = x^2$، المعرَّفة على $\R$:
\begin{enumerate}
  \item \omterm{def:g10:functions:variations}{متناقصة} تمامًا على $\intoc{-\infty}{0}$ \omterm{def:g10:functions:variations}{ومتزايدة} تمامًا
        على $\intco{0}{+\infty}$، وقيمتها الصغرى $0$ عند $x = 0$؛
  \item تحقق $f(-x) = f(x)$ من أجل كل $x$: ومنحنيها، وهو
        \emph{قطع مكافئ}\index{قطع مكافئ}، متناظر بالنسبة إلى المحور
        الشاقولي.
\end{enumerate}
\end{proposition}

\begin{proof}
1. ليكن $0 \leq u < v$. عندئذٍ
\[
f(v) - f(u) = v^2 - u^2 = (v - u)(v + u) > 0,
\]
لأن $v - u > 0$ و $v + u > 0$: إذن $f$ \omterm{def:g10:functions:variations}{متزايدة} تمامًا على
$\intco{0}{+\infty}$. وإذا كان $u < v \leq 0$، فإن $v - u > 0$ لكن
$v + u < 0$، إذن $f(v) - f(u) < 0$: \omterm{def:g10:functions:variations}{متناقصة} تمامًا. وأخيرًا
$x^2 \geq 0 = f(0)$ من أجل كل $x$.

2. لدينا $(-x)^2 = x^2$؛ إذن النقطتان $(x, x^2)$ و $(-x, x^2)$، وهما صورتان
متناظرتان بالنسبة إلى المحور الشاقولي، تقعان معًا على \omterm{def:g10:functions:graph}{المنحنى}.
\end{proof}

\begin{remark}[المربعات والترتيب]
بما أن \omterm{def:g10:functions:function}{دالة} المربع تتناقص على الجهة السالبة، فإن التربيع
\emph{يقلب} ترتيب الأعداد السالبة: $-3 < -2$ لكن $9 > 4$.
لا تربّع طرفي متراجحة أبدًا دون التحقق من الإشارات.
\end{remark}

\section{دالة المقلوب}

\begin{proposition}[دالة المقلوب]\label{prop:g10:reffunc:inverse}
\omterm{def:g10:functions:function}{الدالة} $f(x) = \dfrac1x$، المعرَّفة من أجل $x \neq 0$:
\begin{enumerate}
  \item \omterm{def:g10:functions:variations}{متناقصة} تمامًا على $\intoo{-\infty}{0}$ \omterm{def:g10:functions:variations}{ومتناقصة}
        تمامًا على $\intoo{0}{+\infty}$؛
  \item تحقق $f(-x) = -f(x)$: ومنحنيها، وهو
        \emph{قطع زائد}\index{قطع زائد}، متناظر بالنسبة إلى المبدأ.
\end{enumerate}
\end{proposition}

\begin{proof}
1. ليكن $0 < u < v$. عندئذٍ
\[
f(v) - f(u) = \frac1v - \frac1u = \frac{u - v}{uv} < 0,
\]
لأن $u - v < 0$ و $uv > 0$: إذن \omterm{def:g10:functions:variations}{متناقصة} تمامًا على
$\intoo{0}{+\infty}$. وعلى $\intoo{-\infty}{0}$ يكون للخارج نفسه
$u - v < 0$ و $uv > 0$ من جديد (جداء عددين سالبين)، إذن $f$
\omterm{def:g10:functions:variations}{متناقصة} هناك أيضًا.

2. $\frac{1}{-x} = -\frac1x$.
\end{proof}

\begin{remark}
انتبه: \omterm{def:g10:functions:function}{الدالة} $\frac1x$ \emph{ليست} \omterm{def:g10:functions:variations}{متناقصة} على مجموعة تعريفها كلها. فمن
$u = -1$ إلى $v = 1$ تقفز القيمة من $-1$ صعودًا إلى $1$. ويجب دراسة
الفرعين على حدة.
\end{remark}

\begin{omfigure}
\begin{tikzpicture}
\begin{axis}[omaxis,
  width=5.6cm, height=5.2cm,
  xmin=-3.2, xmax=3.2, ymin=-2.6, ymax=6.4,
  xtick={-2,-1,1,2}, ytick={1,4},
  title={$y = x^2$}]
  \addplot[omDef, thick, domain=-2.45:2.45, samples=90] {x^2};
\end{axis}
\end{tikzpicture}\hfill
\begin{tikzpicture}
\begin{axis}[omaxis,
  width=5.6cm, height=5.2cm,
  xmin=-3.4, xmax=3.4, ymin=-3.4, ymax=3.4,
  xtick={-1,1}, ytick={-1,1},
  title={$y = \frac1x$}]
  \addplot[omDef, thick, domain=0.3:3.1, samples=80] {1/x};
  \addplot[omDef, thick, domain=-3.1:-0.3, samples=80] {1/x};
\end{axis}
\end{tikzpicture}

{\small القطع المكافئ $y = x^2$ (المتناظر بالنسبة إلى المحور الشاقولي) والقطع
الزائد $y = \frac1x$ (فرعان متناظران بالنسبة إلى المبدأ).}
\end{omfigure}

\section{الجذر التربيعي والتكعيب}

\begin{proposition}[دالة الجذر التربيعي]\label{prop:g10:reffunc:sqrt}
\omterm{def:g10:functions:function}{الدالة} $f(x) = \sqrt{x}$، المعرَّفة على $\intco{0}{+\infty}$،
\omterm{def:g10:functions:variations}{متزايدة} تمامًا.
\end{proposition}

\begin{proof}
ليكن $0 \leq u < v$. اضرب واقسم على \emph{المرافق}:
\[
\sqrt v - \sqrt u
= \frac{(\sqrt v - \sqrt u)(\sqrt v + \sqrt u)}{\sqrt v + \sqrt u}
= \frac{v - u}{\sqrt v + \sqrt u} > 0,
\]
لأن $v - u > 0$ و $\sqrt v + \sqrt u > 0$ (ولاحظ أن $v > 0$، إذن
$\sqrt v > 0$).
\end{proof}

\begin{proposition}[دالة التكعيب]\label{prop:g10:reffunc:cube}
\omterm{def:g10:functions:function}{الدالة} $f(x) = x^3$، المعرَّفة على $\R$، \omterm{def:g10:functions:variations}{متزايدة} تمامًا،
ومنحنيها متناظر بالنسبة إلى المبدأ.
\end{proposition}
\admitted

\begin{omfigure}
\begin{tikzpicture}
\begin{axis}[omaxis,
  width=5.6cm, height=5.0cm,
  xmin=-0.7, xmax=5.4, ymin=-0.7, ymax=2.9,
  xtick={1,4}, ytick={1,2},
  title={$y = \sqrt x$}]
  \addplot[omDef, thick, domain=0:5.1, samples=100] {sqrt(x)};
\end{axis}
\end{tikzpicture}\hfill
\begin{tikzpicture}
\begin{axis}[omaxis,
  width=5.6cm, height=5.0cm,
  xmin=-2.2, xmax=2.2, ymin=-3.4, ymax=3.4,
  xtick={-1,1}, ytick={-1,1},
  title={$y = x^3$}]
  \addplot[omDef, thick, domain=-1.48:1.48, samples=90] {x^3};
\end{axis}
\end{tikzpicture}

{\small الجذر التربيعي (المعرَّف من أجل $x \geq 0$ فقط) والتكعيب، وكلاهما
متزايد.}
\end{omfigure}

\begin{proposition}[مقارنة $x$ و $x^2$ و $\sqrt x$]\label{prop:g10:reffunc:compare}
من أجل $0 < x < 1$:\ \ $x^2 < x < \sqrt x$.\quad
ومن أجل $x > 1$:\ \ $\sqrt x < x < x^2$.\quad
وعند $x = 0$ و $x = 1$ تتساوى القيم الثلاث.
\end{proposition}

\begin{proof}
لنفترض أن $0 < x < 1$. بضرب $x < 1$ في $x > 0$ نجد $x^2 < x$.
ثم إن $x < \sqrt x$: فبما أن الطرفين موجبان والتربيع
متزايد على الأعداد الموجبة، فإن هذه المتراجحة تكافئ
$x^2 < x$، وقد برهنّا عليها للتوّ. ومن أجل $x > 1$، يعطي ضرب $x > 1$ في $x$
أن $x^2 > x$، وتنتج $\sqrt x < x$ بحجة التربيع نفسها.
\end{proof}

\begin{omfigure}
\begin{tikzpicture}
\begin{axis}[omaxis,
  width=7.4cm, height=5.8cm,
  xmin=-0.35, xmax=2.5, ymin=-0.35, ymax=2.9,
  xtick={1,2}, ytick={1,2}]
  \addplot[omThm, thick, domain=0:1.65, samples=80] {x^2};
  \addplot[omDef, thick, domain=0:2.35] {x};
  \addplot[omMeth, thick, domain=0:2.35, samples=80] {sqrt(x)};
  \fill (axis cs:1,1) circle (1.5pt);
  \node[omThm, font=\small, right] at (axis cs:1.45,2.4) {$y = x^2$};
  \node[omDef, font=\small, right] at (axis cs:2.0,1.85) {$y = x$};
  \node[omMeth, font=\small, right] at (axis cs:2.0,1.28) {$y = \sqrt x$};
\end{axis}
\end{tikzpicture}

{\small تتقاطع المنحنيات الثلاثة في $(0,0)$ و $(1,1)$ ويتبادل ترتيبها
هناك: فبين $0$ و $1$ يكون $x^2$ الأصغر و $\sqrt x$
الأكبر؛ وبعد $1$ ينقلب الترتيب.}
\end{omfigure}

\begin{method}[مقارنة الصور]\label{met:g10:reffunc:compare}
لمقارنة $f(a)$ و $f(b)$ من أجل \omterm{def:g10:functions:function}{دالة} مرجعية $f$:
\begin{enumerate}
  \item ضع $a$ و $b$ على \omterm{def:g10:numbers:interval}{فترة} تُعرف عليها تغيرات
        $f$؛
  \item إذا كانت $f$ \omterm{def:g10:functions:variations}{متزايدة} هناك، فالصور مرتبة كترتيب
        الدخلين؛ وإذا كانت \omterm{def:g10:functions:variations}{متناقصة}، انقلب الترتيب؛
  \item وإذا لم يكن $a$ و $b$ في \omterm{def:g10:numbers:interval}{فترة} الرتابة نفسها، فعالج
        الإشارات أولًا (مثلًا في حالة المربع: دخل سالب ودخل
        موجب).
\end{enumerate}
\end{method}

\begin{example}
قارن $\dfrac{1}{2.3}$ و $\dfrac{1}{2.4}$: الدخلان كلاهما في
$\intoo{0}{+\infty}$، حيث تتناقص \omterm{def:g10:functions:function}{دالة} المقلوب، و
$2.3 < 2.4$، إذن $\dfrac{1}{2.3} > \dfrac{1}{2.4}$. وقارن $(-1.2)^2$
و $(-1.5)^2$: على $\intoc{-\infty}{0}$ يتناقص المربع، و
$-1.5 < -1.2$، إذن $(-1.5)^2 > (-1.2)^2$ --- وفعلًا $2.25 > 1.44$.
\end{example}

\section{تمارين}

\begin{exercise}[$\star$]\label{exo:g10:reffunc:1}
أعطِ لكل \omterm{def:g10:reffunc:affine}{دالة تآلفية} ميلها وترتيبها عند المبدأ، وقل
هل هي \omterm{def:g10:functions:variations}{متزايدة} أم \omterm{def:g10:functions:variations}{متناقصة}:
\[
f(x) = 3x - 2, \qquad
g(x) = -\tfrac12 x + 4, \qquad
h(x) = 7, \qquad
k(x) = -x .
\]
\end{exercise}

\begin{exercise}[$\star$]\label{exo:g10:reffunc:2}
أوجد \omterm{def:g10:reffunc:affine}{الدالة التآلفية} التي يمر منحنيها بالنقطتين $(2, 1)$ و
$(5, 10)$؛ ثم تلك التي يمر منحنيها بالنقطتين $(-1, 4)$ و $(3, -4)$.
\end{exercise}

\begin{exercise}[$\star$]\label{exo:g10:reffunc:3}
قارن دون آلة حاسبة:
\[
(3.1)^2 \text{ و } (3.2)^2; \qquad
(-2.7)^2 \text{ و } (-2.8)^2; \qquad
\frac{1}{5.1} \text{ و } \frac{1}{5.2}; \qquad
\sqrt{17} \text{ و } \sqrt{15}.
\]
\end{exercise}

\begin{exercise}[$\star$]\label{exo:g10:reffunc:4}
مستعملًا \omterm{def:g10:functions:graph}{منحنى} \omterm{def:g10:functions:function}{دالة} المربع، حل $x^2 = 16$، ثم
$x^2 \leq 16$، ثم $x^2 > 9$.
\end{exercise}

\begin{exercise}[$\star$]\label{exo:g10:reffunc:5}
ليكن $x \in \intoo{0}{1}$. رتّب الأعداد $x$ و $x^2$ و $x^3$ و
$\sqrt x$ من الأصغر إلى الأكبر، وتحقق من جوابك من أجل
$x = 0.25$.
\end{exercise}

\begin{exercise}[$\star\star$]\label{exo:g10:reffunc:6}
حل بيانيًا ثم جبريًا: $\dfrac1x = 2$، ثم
$\dfrac1x < 2$ من أجل $x > 0$. وما الذي يتغير إذا سمحنا أيضًا بالقيم $x < 0$؟
\end{exercise}

\begin{exercise}[$\star\star$]\label{exo:g10:reffunc:7}
علمًا أن \omterm{def:g10:functions:function}{دالة} المربع تتناقص على $\intoc{-\infty}{0}$ وتتزايد
على $\intco{0}{+\infty}$، احصر $x^2$ عندما:
\[
\text{(أ) } x \in \intcc{2}{5};
\qquad
\text{(ب) } x \in \intcc{-3}{-1};
\qquad
\text{(ج) } x \in \intcc{-2}{3}.
\]
(وفي الحالة (ج) انتبه: \omterm{def:g10:functions:extrema}{فالقيمة الصغرى} للمقدار $x^2$ ليست عند طرف.)
\end{exercise}

\begin{exercise}[$\star\star$]\label{exo:g10:reffunc:8}
تتقاضى شركة سيارات أجرة رسمًا ثابتًا قدره $4$ زائد $1.5$ لكل كيلومتر؛
وتتقاضى أخرى $2$ لكل كيلومتر دون رسم. نمذج كل سعر \omterm{def:g10:reffunc:affine}{بدالة
تآلفية} للمسافة $x$، وارسم المستقيمين، وأوجد ابتداءً من أي
مسافة تكون الشركة الأولى أرخص.
\end{exercise}

\begin{exercise}[$\star\star$]\label{exo:g10:reffunc:9}
بيّن أنه من أجل كل $a, b \geq 0$: $\sqrt{a + b} \leq \sqrt a + \sqrt b$.
(إرشاد: الطرفان غير سالبين، فقارن مربعيهما.) ومتى تتحقق
المساواة؟
\end{exercise}

\begin{exercise}[$\star\star\star$]\label{exo:g10:reffunc:10}
لتكن $f(x) = \dfrac{2x + 1}{x - 1}$، معرَّفة من أجل $x \neq 1$.
\begin{enumerate}
  \item بيّن أن $f(x) = 2 + \dfrac{3}{x - 1}$ من أجل كل $x \neq 1$.
  \item استنتج تغيرات $f$ على $\intoo{1}{+\infty}$ من تغيرات
        \omterm{def:g10:functions:function}{دالة} المقلوب.
\end{enumerate}
\end{exercise}

\section{مسألة: قوانين الطبيعة تنطق بالدوال المرجعية}

\begin{problem}[{مسألة نهاية الأسبوع --- مسافات الكبح تنمو مثل
$v^2$، والروافع تخضع للقانون $1/d$، والكواكب تتبع $a^{3/2}$: قراءة
الفيزياء بالدوال المرجعية}]
\label{pb:g10:reffunc:1}
افتح كتاب فيزياء تجد المنحنيات القليلة نفسها في كل صفحة:
القطع المكافئ، والقطع الزائد، \omterm{def:g10:functions:graph}{ومنحنى} الجذر. تكتب الطبيعة
قوانينها بدوال هذا الفصل المرجعية --- ومعرفة
أشكالها (\cref{prop:g10:reffunc:square} و
\cref{prop:g10:reffunc:inverse} و \cref{prop:g10:reffunc:sqrt} و
\cref{prop:g10:reffunc:cube}) تكفي لكبح سيارة، ولموازنة
رافعة، ولضبط توقيت الكواكب. والخاتمة الكبرى هي قانون كبلر
الثالث، مقروءًا مباشرة من معطيات المجموعة الشمسية.

\medskip
\noindent\textbf{الجزء الأول --- قانون الكبح التربيعي.}
قاعدة تقريبية لمسافة كبح سيارة على طريق جاف:
$d = \left(\frac{v}{10}\right)^{2}$ مترًا، عند السرعة $v$ كم/سا.
\begin{enumerate}
  \item احسب مسافات الكبح عند $30$ و $50$ و $90$ و
        $130$ كم/سا.
  \item بكم يضاعف تضعيفُ السرعة مسافةَ الكبح؟
        فسّر ذلك انطلاقًا من تحجيم \omterm{def:g10:functions:function}{دالة} المربع،
        وتحقق منه على $50 \to 100$ كم/سا.
  \item الطب الشرعي: علامات انزلاق تقيس $50$ م. بأي سرعة
        تدين العبارةُ السائقَ (إلى أقرب كم/سا)؟ وأي
        \omterm{def:g10:functions:function}{دالة} مرجعية أجابت --- وعلى أي مجموعة تكون القراءة
        وحيدة (\cref{prop:g10:reffunc:sqrt})؟
  \item قارن المسافة الإضافية الناتجة عن زيادة كم/سا \emph{واحد}،
        عند السرعة المنخفضة وعند السرعة العالية: احسب
        $d(31) - d(30)$ و $d(91) - d(90)$. أي خاصية من خواص
        القطع المكافئ (\cref{prop:g10:reffunc:square}) يوضّحها
        الجوابان؟
  \item يضيف التوقف الحقيقي زمن رد الفعل --- نحو ثانية واحدة،
        تقطع السيارة خلالها $0.28v$ مترًا:
        $D(v) = 0.28v + \left(\frac{v}{10}\right)^2$. احسب
        $D(50)$ و $D(130)$. أي الحدّين --- التآلفي
        أم التربيعي --- يهيمن عند سرعة المدينة، وأيهما على
        الطريق السريع؟
\end{enumerate}

\medskip
\noindent\textbf{الجزء الثاني --- قطوع الحياة اليومية الزائدة.}
\begin{enumerate}[resume]
  \item تتوازن الرافعة عندما تتساوى القوة $\times$ المسافة
        على الجهتين. يجلس طفل كتلته $60$ كغ على بعد $2$ م من
        المحور؛ والقوة الموازِنة على بعد $d$ في الجهة الأخرى
        هي $F(d) = \frac{120}{d}$ (بما يعادل الكيلوغرامات).
        احسب $F(0.5)$ و $F(1)$ و $F(3)$.
  \item ماذا تقول تغيرات \omterm{def:g10:functions:function}{دالة} المقلوب
        (\cref{prop:g10:reffunc:inverse}) عن الروافع ---
        ولماذا يختبئ افتخار أرخميدس (``أعطوني موطئ قدم
        أحرّك لكم الأرض'') في ذيل
        القطع الزائد؟ وما الذي يمنع $d = 0$؟
  \item يخفت الصوت والضوء مع \emph{مربع}
        المسافة: على بعد $1$ م من مصباح تكون الشدة $100$
        وحدة؛ أعطِها على بعد $2$ و $3$ و $10$ م. فسّر
        الأس $2$ بكرة: على أي مساحة انتشر الضوء
        على بعد $d$ (مسألة شاهدة قبر أرخميدس في الكتاب السابق)؟
  \item تكلّف حافلة الرحلة المدرسية $600$ يورو، تُقسم
        بالتساوي على $n$ مشاركًا. اجدول التكلفة على
        الرأس من أجل $n = 10, 20, 30$، وأوجد كم مشاركًا
        ينزل بها إلى $25$ يورو أو أقل، وبيّن ما يعد به شكل
        القطع الزائد --- وما يرفضه --- عندما ينمو $n$.
  \item يكلّف إنتاج $x$ ملصقًا $2x + 3$ يورو (منها $3$ يورو
        تجهيز). بيّن أن التكلفة \emph{المتوسطة} للملصق الواحد هي
        $a(x) = 2 + \frac3x$، وصف تغيراتها،
        وفسّر مقاربها الأفقي اقتصاديًا.
        (قارن بجبر \cref{exo:g10:reffunc:10}.)
\end{enumerate}

\medskip
\noindent\textbf{الجزء الثالث --- انسجام كبلر.}
لكل كوكب، ليكن $a$ بعده المتوسط عن الشمس (بالوحدات
الفلكية، الأرض $= 1$) وليكن $T$ دوره (بالسنوات).
المعطيات: المريخ $a = 1.52$، $T = 1.88$؛ المشتري $a = 5.20$،
$T = 11.86$؛ زحل $a = 9.54$، $T = 29.4$.
\begin{enumerate}[resume]
  \item احسب $a^3$ و $T^2$ من أجل الأرض والمريخ والمشتري.
        ما الذي لاحظه كبلر سنة 1618؟
  \item صُغ القانون على صورة \omterm{def:g10:functions:function}{دالة}: $T = \sqrt{a^3} = a\sqrt a$.
        اختبره على زحل.
  \item تنبؤان: كويكب يدور على $a = 4$ وحدة فلكية ---
        أوجد دوره (تخرج الأعداد صحيحة)؛ ومذنّب
        يعود كل $27$ سنة --- أوجد بعده المتوسط
        (ابحث عن مكعب كامل).
  \item أي ثلاث دوال مرجعية يضفرها القانون
        معًا؟ وما قاعدة تحجيمه: عندما يُضرب $a$
        في $4$، ماذا يحدث للدور $T$؟ (تحقق بالكويكب
        قياسًا على الأرض.)
  \item قرأ كبلر هذا القانون في معطيات تيخو براهي قبل
        سبعين سنة من تفسير جاذبية نيوتن له. في جملة
        واحدة: ماذا تقول هذه الواقعة عن قوة
        التعرّف على \omterm{def:g10:functions:function}{دالة} مرجعية في جدول أعداد؟
\end{enumerate}

\medskip
\noindent\textbf{الجزء الرابع --- السلحفاة والأرنب.}
\begin{enumerate}[resume]
  \item رتّب $\sqrt x$ و $x$ و $x^2$ على $\intoo{0}{1}$ وعلى
        $\intoo{1}{+\infty}$
        (\cref{prop:g10:reffunc:compare})، وتحقق من الترتيبين
        عند $x = 0.25$ وعند $x = 4$.
  \item حل $\sqrt x > x$ حلًّا كاملًا (من أجل $x \geq 0$، ربّع
        بحذر واختم بجدول إشارة).
  \item أدخل التكعيب: رتّب $x$ و $x^2$ و $x^3$ عند
        $x = 0.5$ وعند $x = 2$، وأوجد كل النقط التي يتقاطع
        عندها المربع والتكعيب.
  \item تدفع خطتا ادخار، بعد $t$ سنة، $100\sqrt t$
        يورو (الخطة أ) أو $10 t^2$ يورو (الخطة ب). بيّن أن
        ب تتجاوز أ عندما $t^3 = 100$، وأعطِ لحظة التجاوز
        بالشهر. وما المغزى بشأن الجذور في مواجهة القوى على
        المدى الطويل؟
  \item الخاتمة: من أجل كل قانون في هذه المسألة --- الكبح
        ($v^2$)، والرافعة ($1/d$)، والمصباح ($1/d^2$)، وكبلر
        ($a^{3/2}$)، والخطة أ ($\sqrt t$) --- قل في جملة قصيرة
        أي ملمح من ملامح منحنيه المرجعي يحمل المعنى
        الفيزيائي (الصعود المتزايد الانحدار، والمقارب الشاقولي،
        والكرة المنتشرة، وأس التحجيم، والنمو المتسطّح).
        والخلاصة: الدوال المرجعية هي الأبجدية؛ والطبيعة تكتب بها.
\end{enumerate}
\end{problem}
